\documentclass{amsart}
\usepackage{amsmath}
\usepackage{amssymb}
\usepackage{amsthm}

\newtheorem{theorem}{Theorem}

\title{Right Identity of List Concatenation}
\author{Perqed AI}

\begin{document}

\begin{abstract}
We present a rigorous mathematical proof that appending an empty list to any given list yields the original list. This document translates a machine-verified Lean 4 formalization into a human-readable format.
\end{abstract}

\maketitle

\begin{theorem}
Let $\alpha$ be a type, and let $l$ be a list of elements of type $\alpha$. Then
\[ l \mathbin{++} [] = l \]
where $++$ denotes the list append operation, and $[]$ denotes the empty list.
\end{theorem}

\begin{proof}
We proceed by structural induction on the list $l$. This mirrors the \texttt{induction l} tactic in the formal trace.

\textbf{Base Case:} Suppose $l = []$. We must show that $[] \mathbin{++} [] = []$. This follows immediately from the base definition of the append operation, which states that $[] \mathbin{++} ys = ys$ for any list $ys$. Substituting $ys = []$, the equation holds trivially.

\textbf{Inductive Step:} Suppose $l = x \mathbin{::} xs$ for some element $x \in \alpha$ and some list $xs$. Assume the inductive hypothesis, which states that $xs \mathbin{++} [] = xs$. We wish to show that $(x \mathbin{::} xs) \mathbin{++} [] = x \mathbin{::} xs$. 

By the recursive definition of the append operation, pulling the head element out of the concatenation yields:
\[ (x \mathbin{::} xs) \mathbin{++} [] = x \mathbin{::} (xs \mathbin{++} []) \]
By our inductive hypothesis, we know that $xs \mathbin{++} [] = xs$. Substituting this into our equation, we obtain:
\[ x \mathbin{::} (xs \mathbin{++} []) = x \mathbin{::} xs \]
This establishes the inductive step. 

In the formal Lean 4 trace, the \texttt{simp} tactic chained after the induction (\texttt{<;> simp}) automatically performs these definitional reductions and applies the inductive hypothesis to close both the base and inductive goals. 

Since both the base case and the inductive step hold, the theorem is proven for all lists $l$ by the principle of structural induction.
\end{proof}

\end{document}